\section{共享年代之一：从战后分封到新朝}

这条\hyperref[fig:han-period-timeline]{时间轴}先回答一个定位问题：从前二〇六年的分封到九年的受禅，哪些转换点需要放在同一段长时程中阅读。它只提供年序的定位线，不把并列的事件画成自动因果，也不能替代各地战事、制度和家户经验的不同速度；现在回到\eventref{WH-0206-EIGHTEEN-KINGS}{十八王分封}后的关中，观察第一个转换如何进入道路、仓廪和诸侯关系。
\phantomsection
\label{fig:han-period-timeline}
\begin{center}
\begin{tikzpicture}[x=.035cm,y=1cm,font=\small]
  \fill[paper] (-218,-.6) rectangle (232,2.55);
  \draw[very thick,dynasty] (-206,1.25) -- (220,1.25);

  \foreach \year/\labeltext/\side in {
    -206/{前206\\分封}/above,
    -202/{前202\\汉立}/below,
    -154/{前154\\七国}/above,
    -141/{前141\\武帝}/below,
    -119/{前119\\漠北}/above,
    -91/{前91\\太子危机}/below,
    -9/{前9\\王莽居摄}/above,
    9/{9\\新}/below,
    23/{23\\更始}/above,
    25/{25\\复汉}/below,
    89/{89\\北边纪功}/above,
    184/{184\\黄巾}/below,
    189/{189\\宫门失守}/above,
    200/{200\\官渡}/below,
    208/{208\\赤壁}/above,
    220/{220\\禅代}/below}
  {
    \draw[timeline,thick] (\year,1.12) -- (\year,1.38);
    \node[\side,align=center,font=\scriptsize,text=ink] at (\year,1.42) {\labeltext};
  }

  \draw[dashed,source!75] (9,.08) -- (9,2.3);
  \draw[dashed,source!75] (25,.08) -- (25,2.3);
  \node[align=center,font=\scriptsize,text=dynasty] at (-100,.15) {西汉：关中、郡国、边疆与财政的重组};
  \node[align=center,font=\scriptsize,text=timeline] at (17,2.27) {新与复汉：命令、名分和执行的断裂};
  \node[align=center,font=\scriptsize,text=dynasty] at (124,.15) {东汉：雒阳复接、地方军权与区域竞争};

  \node[draw=source!75,fill=paper,rounded corners=2pt,align=left,
    inner sep=3pt,font=\scriptsize] at (105,2.45)
    {时间轴为编年阅读的定位线，不表示事件之间\\
     存在自动因果，也不按战争规模或政权控制范围绘制。};
\end{tikzpicture}

\smallskip
\small\textit{图\quad 汉帝国时期时间轴（前206--220）：据《史记》《汉书》《后汉书》《后汉纪》《三国志》与《资治通鉴》的可核年序编制。它只标出本卷反复回看的转换点；政权更替、战争与制度动作的精确日期、参与者和影响范围仍以各节的来源注记为准。\quad\hyperref[sec:han-shared-chronology]{回到共享年序}}
\end{center}


前二〇六年的关中有仓廪、关隘、县廷和道路，却没有一位能令所有军队听命的主人。项羽在十八王分封中重排秦故地，自立西楚霸王，刘邦被徙往汉中。封号暂时安排军功，不能决定旧吏、粮食和部众此后归谁；鸿门会的行动次序可以大体安放，人物言辞和谋划则主要来自《史记》的后出叙事，不能当作逐时记录。

刘邦还定三秦、东向争夺的过程，须放在多个后方一起看。彭城溃败后的重组，使荥阳、成皋、敖仓、关中和河北的道路、兵员、粮赋同时成为胜负条件。到垓下合围，项羽的败亡来自联盟、供给与战场退路的共同压力；楚歌和人物绝唱属于最有力的文学化叙事，不能代替这些条件。前二〇二年二月，刘邦即皇帝位，把战时联盟置于汉的名号之下，却没有消除功臣、诸侯王和地方网络的独立分量。

白登之围使新朝很早看见北边的尺度。前二〇〇年冬，刘邦北征受困；和亲、关市和边郡守备并非“强”或“弱”的单一选择，而是在骑兵、道路与财力有限时争取恢复时间的组合。此后处置异姓王、以刘氏宗室置换封国，令中枢较能控制可独立动员的将领，却也保留了郡国并行的长期张力。

文、景时期的轻徭、刑制调整、田租政策和铸钱取向，不能由后世“文景之治”的概括替代。诏令最能说明朝廷希望如何恢复编户和生产，难以逐户证明赋役与土地占有。张力在前一五四年公开化：削藩与七国举兵以削地为触发因素，吴楚等王能够起兵还因其财富、封国官属与地方关系；周亚夫在昌邑断粮道，才令中央取得胜势。战后限制王国任官和司法，扩大了朝廷核验地方军政的能力，并未抹掉地方家族、市场与城邑。

武帝朝把这套能力推向边界。前一三九年张骞出使带来的首先是片段信息，而非一张准确的西域地图。马邑伏击失败后，汉匈关系转入大战；卫青经营朔方、霍去病在河西作战和前一一九年的漠北出击改变了战略条件。战绩、降附部众和设郡记录了朝廷的军事与行政动作，不能据此把绿洲、牧地或走廊写成即刻而均匀的控制区。

南越之战也有自己的地理。前一一二年吕嘉杀汉使后，朝廷发兵，次年入番禺并置郡；海路、山地、地方首领和迁入者决定新任命能走多远。远征要求把钱币、盐铁、转运与军费接进同一张网。前一一九年的五铢与盐铁措施、均输和平准的相关记录，说明中枢试图取得更多财货与信息；《食货志》的制度叙述和出土钱币各有尺度，不能被合并成全国收支表。

前九一年太子危机使继承和都城军队相撞；人物告发与巫蛊细节多需从后出的叙事辨认。前八九年轮台诏拒绝继续大规模屯田远征，至少表明朝廷公开衡量远距供给的代价，不等于边政从此退回关中。武帝死后，霍光等辅政，前七四年昌邑王被废、宣帝被立。废立事实和年月可核，失德清单是为政治行动辩护的文本，不能逐项复原宫廷日常。

宣帝朝的赵充国在羌地争论中主张屯田、招抚与分化，显示边政并非只有进攻一途；“羌”是战争和行政材料中的分类，不能替代各部众与边民的多样关系。前五十一年呼韩邪入朝，朝廷以印绶和赏赐经营北边关系，同样不是一纸称臣可以穷尽的服属。

元、成、哀、平之际，王氏逐渐占据大司马和人事枢纽。前八年王莽任大司马，前一年一度退居，说明外戚上升并非无摩擦的直线。元始二年的户口统计是朝廷簿籍的行政申报，能够证明国家仍以户口表达治理尺度，不能直接等同真实总人口、每户负担或地方执行率。平帝死后，居摄、符命与劝进逐步排空孺子婴；九年新朝建立，把已取得的军政、礼仪和文书位置转换为受禅的仪式。

\begin{turningpoint}[前二〇六年至九年留下的压力]
战后封国、削藩、扩张、财政集中与辅政都解决了眼前的协调难题：谁供粮，谁任官，谁守住道路，谁能在幼主时传递命令。每一次解决也制造新的接口和新的争夺者。新朝并不是西汉行政网络突然消失后的空白；它显示名分、符命和改革语言可以被掌握中枢接口的人使用，而这些语言能否越过乡里、市场和边郡，仍须接受执行的检验。
\end{turningpoint}
